\chapteropen
\chapter{UCD formation channels}
\label{app:ucd_channels}

The consensus of current research on possible formation processes for \glspl{ucd} \citep[e.g.][]{Norris_2011,Pfeffer_2013,Penny_2014,Norris_2014,Pfeffer_2014,Norris_2015,Pfeffer_2016,Goodman_2018,Mahani_2021,Saifollahi_2021,Mayes_2021, Pomeroy_2025} suggests multiple formation paths. \Glspl{ucd} may represent a continuation of the high-mass end of the \gls{gc} population formed through normal massive star-cluster formation. Alternatively they could be the product of tidal stripping after a nucleated dwarf galaxy encounters a larger galaxy, removing the majority of the dwarf's extended outer structure, while leaving the core relatively intact \citep{Bekki_2001}. Multiple formation channels and thus a composite origin is widely supported, with overlap below a star cluster formation limit of $M_\star\lesssim\!5\times10^7\Msun$ \citep{Norris_2019}. Another plausible alternative for the formation of these objects is mergers of multiple star clusters which have formed in cluster complexes, often after gas-rich mergers and starburst episodes \citep{Fellhauer_2002, Bruns_2011}.

Different formation channels can produce structurally similar present-day systems, therefore, observational discriminants are essential. Evidence for a central black hole, complex star-formation histories, elevated metallicity at fixed mass, internal abundance spreads, and size--mass scaling relations can all help distinguish stripped nuclei from GC-like \glspl{ucd} \citep[\eg][]{Pfeffer_2016,Norris_2019}. In general, stripped nuclei are expected to appear more chemically complex and more metal-rich for their mass than systems formed primarily as massive star clusters.

Because these different channels can produce structurally similar present-day systems, the following observational discriminants are commonly used to distinguish candidate stripped nuclei from GC-like \glspl{ucd}.

\begin{itemize}
    \item Dwarf ellipticals likely contain \gls{smbh}, and in their privileged central nuclear position are unlikely to be affected by any stripping activity. Thus the presence of a central \gls{smbh} or \gls{imbh}, observable through an elevated dynamical mass-light ratio, would identify the \gls{ucd} as a stripped nucleus \citep{Pfeffer_2016}.
    \item Similarly, the nucleus of a progenitor dwarf galaxy is likely to have been through multiple episodes of star formation, and thus evidence of a complex star formation history is likely indicative a \gls{ucd} is a remnant nucleus. This view is, however, nuanced as a \gls{ucd} formed from multiple star clusters may also exhibit a complex star formation signature.
    \item Patterns of metallicity, stellar abundance and size--mass trends can also indicate whether an object sits on a GC-like sequence or looks more like nuclear star clusters. For example, GC-like objects have been shown to follow a mass-metallicity relation, also known as a \gls{blue-tilt}, where higher masses of cluster show increasing metallicity. This is interpreted as self-enrichment during formation, where higher-mass clusters have a greater gravitational well to retain gas, dust and stellar processing remnants. Conversely, stripped nuclei are often too metal-rich for their mass compared to \glspl{gc}, as they follow galaxy mass-metallicity relations, where the metallicity reflects the host galaxy potential well (not cluster self-enrichment). 
    \item As summarized in \cref{tab:gc_vs_ucd_spectra}, other metallicity and abundance trends include internal metallicity spreads, where stripped nuclei are likely to show measurable [\ce{Fe}/\ce{H}] spreads, compared to the more homogeneous \glspl{gc}. Also, \gls{alpha-element} abundance is typically enhanced in \glspl{gc} due to rapid star formation dominated by core-collapse \glspl{sn}, compared to the longer formation timescales of stripped nuclei, which also have more contribution from Type 1a \glspl{sn}. Finally, the age distribution of GC-like \glspl{ucd} are likely narrow and old ($\gtrsim\!\qtyrange{10}{12}{\giga\year}$) compared to the composite ages of stars in stripped nuclei which may have an intermediate and young stellar component, in addition to an old base population.
    
\end{itemize}

\begin{table}[htbp]
\centering
\caption[Spectral Differences between GCs and UCDs.]{Spectral Differences between \glspl{gc} and \glspl{ucd}. Details from high signal-to-noise spectra can uncover subtle structural and chemical differentiators between \glspl{gc} and \glspl{ucd}}
\label{tab:gc_vs_ucd_spectra}
\begin{adjustbox}{max width=\textwidth}
\begin{tabular}{p{2.5cm} p{3.8cm} p{4cm} p{4.2cm}}
\doubletoprule
\textbf{Spectral Feature} & \textbf{Globular Clusters} & \textbf{Ultracompact Dwarfs} & \textbf{Diagnostic Role in Classification} \\ 
\midrule
Line Width / Velocity Dispersion ($\sigma$) & Narrow absorption features; $\sigma$ typically under 20 km/s. & Broader absorption features; $\sigma$ generally ranges between 20 and 50 km/s. & Broader lines indicate higher mass and a deeper gravitational well, typical of \glspl{ucd}. \\ \hline
Metallicity \& Abundances & Highly ancient, metal-poor to intermediate profiles ([Fe/H] up to -0.5). & Generally higher global metallicities, often reaching solar or near-solar levels. & Higher metallicity points toward complex chemical enrichment or a stripped galactic nucleus. \\ \hline
\gls{alpha-abundance} ([$\alpha$/Fe]) & Uniformly elevated [$\alpha$/Fe] ratios (roughly +0.3 to +0.4 dex). & Highly variable [$\alpha$/Fe] ratios; can range from solar up to highly alpha-enhanced. & Uniform enhancement indicates a single, rapid burst of star formation early in cosmic history. \\ \hline
Balmer Absorption Lines (H$\beta$, H$\gamma$, H$\delta$) & Sharp, weak to moderate hydrogen lines consistent with old stellar turn-off ages. & Occasionally show stronger Balmer lines, indicating younger or intermediate-age sub-populations. & Traces the overall age distribution and reveals if the object experienced prolonged star formation. \\ \hline
Continuum Shape (UV/Optical) & Stark, older stellar continuum peaking in the red/near-IR wavelengths. & Occasionally shows blue horizontal branch or young stellar continuum excesses. & Confirms the lack of ongoing star formation and helps estimate overall integrated spectral type. \\ \hline
Central Black Hole Signatures & No broad emission or dynamic line-splitting at center. & Can show subtle nuclear line-splitting or non-thermal continuum components. & Confirms the presence of a supermassive black hole, verifying it as a stripped dwarf galaxy core. \\ \hline
\end{tabular}
\end{adjustbox}
\tablecomments{Summarized from 
\citet{Drinkwater_2003, Evstigneeva_2007, Seth_2014}
}
\end{table}

